\documentclass[11pt,a4paper]{article}

\usepackage[T1]{fontenc}
\usepackage[utf8]{inputenc}
\usepackage[french]{babel}
\usepackage{lmodern}
\usepackage[margin=2.5cm]{geometry}
\usepackage{amsmath,amssymb}
\usepackage{graphicx}
\usepackage{hyperref}

\title{Titre de l'article}
\author{Prénom Nom\\ \small Établissement, Ville, Pays\\ \small \texttt{adresse@exemple.org}}
\date{\today}

\begin{document}

\maketitle

\begin{abstract}
Résumé de l'article en quelques lignes : contexte, objectif, méthode, principaux résultats.
\end{abstract}

\section{Introduction}
Présentez le contexte et la question étudiée. Les références se citent ainsi~\cite{knuth1984}.

\section{Méthode}
Décrivez la démarche. Une équation numérotée :
\begin{equation}
  E = mc^2 \label{eq:energie}
\end{equation}
L'équation~\eqref{eq:energie} se référence automatiquement.

\section{Résultats}
\begin{table}[h]
  \centering
  \begin{tabular}{lcc}
    \hline
    Échantillon & Valeur A & Valeur B \\
    \hline
    1 & 12 & 15 \\
    2 & 14 & 18 \\
    \hline
  \end{tabular}
  \caption{Exemple de tableau.}
  \label{tab:resultats}
\end{table}

\section{Conclusion}
Résumez les apports et les perspectives.

\bibliographystyle{plain}
\bibliography{references}

\end{document}
